% ---------------------------------------------------------------------
\section{Problem Formulation}\label{sec:problem}
% ---------------------------------------------------------------------

\subsection{Station, time and prices}\label{sec:station}

A CPO runs a station with $N$ chargers $\mathcal N=\{1,\dots,N\}$ behind one transformer, and each charger hosts at most one EV.
While an EV is connected, its quantities carry the index $i$ of its charger.
Time is divided into steps $t=0,\dots,T-1$ of length $\Delta t$ [h], and the action of step $t$ is held constant on $[t\Delta t,(t+1)\Delta t)$.
At every step the station observes the transformer rating $P_{\sm{tr}}(t)$, the inflexible load $P^{\sm{inf}}(t)\ge0$ and the photovoltaic (PV) generation $P^{\sm{pv}}(t)\ge0$, all in kW.
Energy is bought at the price $c^{\sm{ch}}(t)$ and sold at the price $c^{\sm{dis}}(t)$ [\texteuro/kWh], and both prices may be negative.

\subsection{EV energy dynamics}\label{sec:dynamics}

EV $i$ arrives at step $t_i^{\sm{arr}}$ and is controlled at the steps $t=t_i^{\sm{arr}},\dots,d_i$, where the departure step $d_i$ is the last controlled step.
The EV leaves after the action of step $d_i$ has been applied, so with $E_i(t)$ denoting the battery energy at the start of step $t$, the departure energy is $E_i(d_i+1)$.
At step $t$, the number of remaining actions, including the current one, is
\begin{equation}\label{eq:tau}
\tau_i(t)=d_i-t+1 ,
\end{equation}
and the set of occupied chargers is $\mathcal N(t)=\{i: t_i^{\sm{arr}}\le t\le d_i\}$, with $n(t)=|\mathcal N(t)|$ connected EVs.

Let $p_i(t)$ [kW] be the grid-side power of charger $i$, positive for charging and negative for V2G discharging.
With $p^{+}=\max(p,0)$ and $p^{-}=\max(-p,0)$, the battery energy evolves as
\begin{equation}\label{eq:dynamics}
E_i(t+1)=E_i(t)+\Delta t\,g_i\big(p_i(t)\big),\qquad g_i(p)=\eta_i^{\sm c}p^{+}-p^{-}/\eta_i^{\sm d},
\end{equation}
where $\eta_i^{\sm c},\eta_i^{\sm d}\in(0,1]$ are the charge and discharge efficiencies, and the energy map $g_i$ is continuous, strictly increasing and piecewise linear with $g_i(0)=0$.
Per-step physical limits restrict the power to
\begin{equation}\label{eq:limits}
\pmn_i(E)\le p_i(t)\le \bar P_i(E),\qquad \pmn_i(E)=-\min\Big\{\Pdis_i,\ \eta_i^{\sm d}\frac{(E-E_i^{\min})^{+}}{\Delta t}\Big\},\qquad E=E_i(t),
\end{equation}
where $\bar P_i(E)\ge0$ is the state-dependent maximum charging power.
This maximum includes the constant-current, constant-voltage taper and the capacity limit $\eta_i^{\sm c}\bar P_i(E)\Delta t\le E_i^{\sm{cap}}-E$, and its nominal value is $\Pnom_i=\sup_E\bar P_i(E)$.
The maximum discharging power is $\Pdis_i\ge0$, and under \Eq{eq:limits} the energy stays in $[E_i^{\min},E_i^{\sm{cap}}]$.

The theory of \Sec{sec:method} holds for the general map $g_i$, but EV2Gym reports, bills and meters the battery-side energy change as the charger power \cite{orfanoudakis2025ev2gym}.
Our experiments therefore use $\eta_i^{\sm c}=\eta_i^{\sm d}=1$, so that $g_i(p)=p$ and EV2Gym's own efficiencies act only as derating factors on the power limits.
This convention removes the round-trip loss from V2G arbitrage, and \Sec{sec:limitations} lists it as a limitation.

\subsection{Departure energy requirement}\label{sec:requirement}

Each EV requests a desired departure energy $E_i^{\sm{des}}$, and at arrival the request is capped to the energy that is individually reachable during the stay,
\begin{equation}\label{eq:cap}
\Ereq_i=\min\Big\{E_i^{\sm{des}},\ E_i^{\sm{arr}}+R_i\big(\tau_i(t_i^{\sm{arr}}),E_i^{\sm{arr}}\big)\Big\},
\end{equation}
where $E_i^{\sm{arr}}=E_i(t_i^{\sm{arr}})$ and $R_i(k,E)$ is the reachable-energy operator of \Sec{sec:bounds}.
Without this cap, some arrivals could not be served by any controller, and a guarantee about them would be void.
The departure energy requirement of EV $i$ is therefore
\begin{equation}\label{eq:requirement}
E_i(d_i+1)\ \ge\ \Ereq_i .
\end{equation}
In code, the cap uses the same certified threshold routine as the safety layer and lowers the value of \Eq{eq:cap} by at most $10^{-6}$ kWh.
Admission control (\Sec{sec:fallback}) may lower the requirement of a newly arriving EV further, and the experiments report the number and size of capped and reduced requests.

\subsection{Transformer limit}\label{sec:transformer}

The net transformer power $P^{\sm{inf}}(t)-P^{\sm{pv}}(t)+\sum_{i}p_i(t)$ must stay within $[-P_{\sm{tr}}(t),P_{\sm{tr}}(t)]$, and for the EV aggregate this symmetric thermal limit gives the transformer limit
\begin{equation}\label{eq:transformer}
-\Crev(t)\le\sum_{i\in\mathcal N}p_i(t)\le C(t),\qquad
C(t)=P_{\sm{tr}}(t)-P^{\sm{inf}}(t)+P^{\sm{pv}}(t),\qquad \Crev(t)=P_{\sm{tr}}(t)+P^{\sm{inf}}(t)-P^{\sm{pv}}(t).
\end{equation}
The forward limit $C(t)$ bounds charging and the reverse limit $\Crev(t)$ bounds V2G export, and their sum $2P_{\sm{tr}}(t)$ is never negative.
Empty chargers draw zero power, and the safety layer uses the measured inflexible load and PV of the current step.
The policy, in contrast, observes only noisy forecasts of future load and PV.

\subsection{Online control problem}\label{sec:online}

At each step the CPO chooses $p(t)\in\R^N$ using only the information available at step $t$ and no knowledge of future arrivals.
This information comprises the arrival data and current energies of the connected EVs, current and short-window day-ahead prices, and the measured load and PV with their forecasts.
The objective is the expected episode profit
\begin{equation}\label{eq:objective}
\max\ \mathbb E\Big[\sum_{t=0}^{T-1}\Delta t\sum_{i\in\mathcal N(t)}\big(c^{\sm{dis}}(t)\,p_i^{-}(t)-c^{\sm{ch}}(t)\,p_i^{+}(t)\big)\Big],
\end{equation}
subject to \Eqs{eq:dynamics}{eq:limits} at every step, \Eq{eq:requirement} for every EV and \Eq{eq:transformer} at every step.

\subsection{Perfect-information optimum}\label{sec:oracle}

An oracle that knows the whole episode provides the reference for the optimality gap.
It sees all sessions with $(t_i^{\sm{arr}},d_i,E_i^{\sm{arr}},\Ereq_i)$, all prices and the realised limits $C(t)$ and $\Crev(t)$.
For every session and every step of its stay, the oracle uses charging and discharging powers $p_i^{\pm}(t)\ge0$, energies $E_i(t)$ and a shortfall slack $z_i\ge0$ under the constraints
\begin{equation}\label{eq:lp}
\begin{aligned}
&E_i(t_i^{\sm{arr}})=E_i^{\sm{arr}},\qquad E_i(t+1)=E_i(t)+\Delta t\big(\eta_i^{\sm c}p_i^{+}(t)-p_i^{-}(t)/\eta_i^{\sm d}\big),\\
&E_i^{\min}\le E_i(t)\le E_i^{\sm{cap}},\qquad p_i^{+}(t)\le\Pnom_i,\qquad p_i^{-}(t)\le\Pdis_i,\qquad E_i(d_i+1)+z_i\ge\Ereq_i,\\
&-\Crev(t)\le\textstyle\sum_{i}\big(p_i^{+}(t)-p_i^{-}(t)\big)\le C(t).
\end{aligned}
\end{equation}
The linear program (LP) is solved lexicographically with the HiGHS solver \cite{huangfu2018parallelizing}: stage 1 minimises the total shortfall $\sum_iz_i$, and its optimum $z^{*}$ is zero whenever all requirements are jointly feasible.
Stage 2 then maximises the profit of \Eq{eq:objective} subject to \Eq{eq:lp} and $\sum_iz_i\le z^{*}+10^{-6}$, and its optimum is $J^{\sm{LP}}$.

The LP relaxes the state-dependent taper to the nominal power $\Pnom_i$ and allows simultaneous charging and discharging, which never pays when $c^{\sm{dis}}\le c^{\sm{ch}}$ and is therefore exact in our scenarios.
In the exact model, $J^{\sm{LP}}$ bounds the profit $J^{\pi}$ of every controller $\pi$ that respects \Eq{eq:transformer} and leaves a total shortfall of at most $z^{*}$.
EV2Gym, however, rounds the stored energy of an active EV up to 0.01 kWh after every step and does not bill this energy.
A feasible controller inside the simulator can therefore meet its targets while buying slightly less energy than the exact model requires, and its profit may exceed $J^{\sm{LP}}$.
We therefore use the rounding-aware optimum $J^{\sm{LP,ceil}}$ of the same LP with $\Ereq_i$ and $E_i^{\min}$ lowered by the largest possible rounding gain of 0.01 kWh per step of the stay.
Proposition~\ref{S-prop:oracle} in the supplementary material shows that $J^{\sm{LP,ceil}}$ bounds the profit of every such controller inside EV2Gym, and the optimality gap is defined as
\begin{equation}\label{eq:gap}
\mathrm{gap}(\pi)=\frac{J^{\sm{LP,ceil}}-J^{\pi}}{|J^{\sm{LP,ceil}}|}.
\end{equation}
Profits are usually negative, because charging to the requested energy costs money, so the absolute value is needed in the denominator.
Episodes with $|J^{\sm{LP,ceil}}|<1$\,\texteuro{} are excluded from this relative metric, and the gap with respect to $J^{\sm{LP}}$ is reported in the supplementary material.
Because of the relaxations, the reported gap overestimates the true gap; controllers that violate the transformer limit are not covered by the bound and are flagged in the results.

\subsection{Markov decision process}\label{sec:mdp}

We cast the online problem as an MDP with one decision per step, although the problem is strictly partially observable because future arrivals, prices beyond a short window and future loads are not observed.
Following standard practice, the policy conditions on the observation described in \Sec{sec:policy}.

\textbf{Action.} The policy outputs one raw action $a_i(t)\in\R$ per charger, and the action is scaled to kW as $x_i(t)=P^{\sm{ref}}a_i(t)$ with $P^{\sm{ref}}=11$ kW.
Its executed counterpart $p(t)=\Pi_t(x(t))$ is the output of the safety layer of \Sec{sec:projection}, or of the fallback rule of \Sec{sec:fallback} when the feasible set is empty.
No squashing function is used, so all bounding is done by $\Pi_t$.

\textbf{Reward.} The reward is the realised profit of the step,
\begin{equation}\label{eq:reward}
r(t)=\Delta t\sum_{i\in\mathcal N(t)}\big(c^{\sm{dis}}(t)\,\tilde p_i^{-}(t)-c^{\sm{ch}}(t)\,\tilde p_i^{+}(t)\big),
\end{equation}
where $\tilde p_i$ is the power reported by the simulator.
The reward contains no penalty for the departure energy requirement or the transformer limit, because the safety layer enforces both and such a penalty would vanish on every step where aggregate feasibility holds.
On infeasible steps the fallback rule is deterministic, so a penalty there would add variance without a learnable signal.
For training, the reward is divided by $N r^{\sm{ref}}$ with $r^{\sm{ref}}=c^{\sm{ref}}P^{\sm{ref}}\Delta t$ and $c^{\sm{ref}}=0.25$\,\texteuro/kWh, so that returns do not grow with $N$ and one critic can serve stations of different sizes.
All reported results, in contrast, use the undivided profit $r(t)$ in \texteuro.

\textbf{Transition and horizon.} Arrivals, prices, loads and PV evolve independently of the actions, and energies follow \Eq{eq:dynamics}.
An arriving EV adds a token and an EV with $d_i=t$ leaves after step $t$, and since every EV departs before the episode ends, step $T$ is a true terminal state.
The discount factor $\gamma=0.99$ gives an effective horizon of about one day at $\Delta t=15$ min.
